\section{Stability of the comparison grid}
\label{sec:stability}

The comparison grid supplies a reference colouring without prescribing any
target edge colours.  We first prove a quantitative continuous statement,
then transfer it to every finite rectangular grid.  The finite statement
applies to minimizers with arbitrary bounded unary perturbations.
A separate estimate applies to any colouring with the resulting localized
support; this distinction will allow us to change the program colours while
keeping the same grid colouring.

\subsection{The continuous comparison energy}

All integrals in this section use Lebesgue measure.  Functions are Borel
measurable; the statements also hold for Lebesgue-measurable functions after
replacement on a null set.  For a sign function
$f:[0,1]^2\to\{-1,1\}$, put
\begin{equation}
 \mathcal E(f)=
 \int_{\substack{0\le u<u'\le1\\0\le v<v'\le1}}
 \mathbf 1\{f(u,v)=f(u',v')\}\,du\,dv\,du'\,dv'.
 \label{eq:comparison-energy}
\end{equation}
Each comparable unordered pair occurs once.  Define
\[
 f_*(u,v)=
 \begin{cases}1,&u+v<1,\\-1,&u+v\ge1.\end{cases}
\]
The diagonal convention has no effect on the continuous integral, but we
retain this strict inequality when passing to finite grids.

\begin{theorem}[Continuous two-phase stability]
\label{thm:continuous}
Every measurable sign function on the unit square has
$\mathcal E(f)\ge1/12$.  For every $\varepsilon\ge0$, if
$\mathcal E(f)\le1/12+\varepsilon$, then
\begin{equation}
 \min_{\sigma\in\{-1,1\}}
 \|f-\sigma f_*\|_{L^1([0,1]^2)}
 \le4096\varepsilon^{1/72}.
 \label{eq:continuous-stability}
\end{equation}
\end{theorem}

We give the analytic ingredients with constants, since only a uniform
estimate, rather than a sharp exponent, is needed later.

For $h:[0,1]\to\{-1,1\}$ define
\[
 (\mathcal T h)(t)=\int_0^1\operatorname{sgn}(s-t)h(s)\,ds,
 \qquad D(h)=\frac13-\|\mathcal T h\|_2^2.
\]
Values of the sign kernel on its diagonal are immaterial.  A threshold
function is a sign function that is constant on each side of a point
$c\in[0,1]$, with opposite signs on those two intervals.  Both orientations
and, through $c=0,1$, both constant functions are included.

\begin{lemma}[One-dimensional stability]
\label{lem:one-dimensional}
For every measurable sign function $h$, one has $D(h)\ge0$.  There is a
threshold function $b$ such that
\[
 \|h-b\|_1\le4D(h)^{1/3}.
\]
\end{lemma}

\begin{proof}
Let
\[
 I(h)=\int_{0<r<s<t<1}
 \mathbf 1\{h(r)=h(t)\ne h(s)\}\,dr\,ds\,dt.
\]
The elementary identity
\[
 4\mathbf 1\{h(r)=h(t)\ne h(s)\}
 =1+h(r)h(t)-h(r)h(s)-h(s)h(t)
\]
gives
\begin{equation}
 D(h)=8I(h).
 \label{eq:alternation-deficit}
\end{equation}
To check the coefficient, after integrating out the third variable the
coefficient of $h(u)h(v)$, for $u<v$, in the triple expression is
$2(v-u)-1$.  On the other hand,
\[
 \int_0^1\operatorname{sgn}(u-t)\operatorname{sgn}(v-t)\,dt
 =1-2|u-v|.
\]
The constant triple integral is $1/6$, and the ordered two-variable
integral is half the corresponding symmetric integral.  These observations
give \eqref{eq:alternation-deficit}, including its constant term, by Fubini.
All integrands are bounded, so these changes of order are legitimate.

Let $e$ be the least measure of disagreement with a threshold.  For each
orientation, disagreement is continuous in its threshold position, and
$[0,1]$ is compact, so the minimum is attained.  We claim
$I(h)\ge e^3/64$.  There is nothing to prove if $e=0$.  Otherwise let $p$
be the measure of the positive set.  The constant negative threshold shows
that $p\ge e$.  By continuity of cumulative positive mass choose $a\le b$
so that the positive mass before $a$ and after $b$ is $e/4$ each.

If the negative mass between $a$ and $b$ is at least $e/4$, independently
choosing a positive point before $a$, a negative point between $a,b$, and
a positive point after $b$ gives $I(h)\ge(e/4)^3$.  Otherwise compare with
the threshold positive before $b$ and negative after $b$.  Its disagreement
is at least $e$, whereas the negative mass between $a,b$ plus the positive
mass after $b$ is less than $e/2$.  Thus the negative mass before $a$ is
at least $e/2$.  Comparing with the oppositely oriented threshold at $a$
similarly gives negative mass at least $e/2$ after $b$.  The positive mass
between $a,b$ is $p-e/2\ge e/2$.  Negative--positive--negative triples
therefore give the stronger bound $I(h)\ge e^3/8$.  The claim follows.
By \eqref{eq:alternation-deficit}, $e^3\le8D(h)$; since the sign-function
$L^1$ distance is $2e$, the asserted bound follows.
\end{proof}

Apply $\mathcal T$ in the two coordinates and write
\[
 A(u,v)=\int_0^1\operatorname{sgn}(s-u)f(s,v)\,ds,
 \qquad
 C(u,v)=\int_0^1\operatorname{sgn}(s-v)f(u,s)\,ds,
 \qquad M=\int_{[0,1]^2}f.
\]
Expanding the order indicators and the equality indicator, and using
antisymmetry of the sign kernel, gives
\[
 \mathcal E(f)=\frac{1+M^2-\langle A,C\rangle}{8}.
\]
In particular,
\begin{equation}
 8\bigl(\mathcal E(f)-1/12\bigr)
 =M^2+\frac12\|A-C\|_2^2
 +\frac12\left(\frac13-\|A\|_2^2\right)
 +\frac12\left(\frac13-\|C\|_2^2\right).
 \label{eq:comparison-deficit}
\end{equation}
Lemma~\ref{lem:one-dimensional} on each row and column makes every term on
the right nonnegative.  This proves the lower bound in
Theorem~\ref{thm:continuous}.  For $f_*$ the mean vanishes, both slice
families are thresholds, and direct integration gives $A=C$.  Hence its
energy is exactly $1/12$.

Under an excess bound $\varepsilon$, \eqref{eq:comparison-deficit} gives
\begin{equation}
 |M|\le\sqrt{8\varepsilon},\qquad
 \|A-C\|_2\le4\sqrt\varepsilon,
 \qquad
 \int_0^1 D(f(\cdot,v))\,dv\le16\varepsilon,
 \label{eq:deficit-consequences}
\end{equation}
and the analogous column estimate.  Jensen's inequality and
Lemma~\ref{lem:one-dimensional} show that the average row distance from a
threshold, and also the average column distance, is at most
$4(16\varepsilon)^{1/3}\le11\varepsilon^{1/3}$.  These distance functions are
measurable: one may take the infimum over rational threshold positions in
both orientations, using continuity in the threshold position.  No
measurable selection of optimal thresholds is required.

\begin{lemma}[Approximation by a diagonal profile]
\label{lem:diagonal-approximation}
Suppose $0<h\le1$ and $\mathcal E(f)\le1/12+h^6$.  There is a measurable
$b:[0,2]\to\{-1,1\}$ such that, for $f_1(u,v)=b(u+v)$,
\[
 \|f-f_1\|_1\le128h,
 \qquad \mathcal E(f_1)\le1/12+257h.
\]
\end{lemma}

\begin{proof}
It suffices first to treat $h<1/8$.  Extend $f$ by zero outside the square.
A threshold on $[0,1]$, extended by zero, has total variation at most four.
Approximating each row by a threshold and applying the triangle inequality
shows that translation by $r$ in the first coordinate changes the zero
extension in $L^1$ by at most
$22\varepsilon^{1/3}+4|r|$.  The same holds in the second coordinate.
Here the rowwise inequality can be integrated using the measurable
infimum just described, so no selection issue arises.

Let $f_h$ be the convolution of the zero extension with the uniform
probability density on $[-h,h]^2$.  Integrating the translation bound and
using $\varepsilon=h^6$ yields
\begin{equation}
 \|f-f_h\|_{L^1(\mathbb R^2)}
 \le44h^2+8h\le52h.
 \label{eq:box-error}
\end{equation}
On the interior square $Q_h=[h,1-h]^2$ we also have
\begin{equation}
 \| (\partial_u-\partial_v)f_h\|_{L^1(Q_h)}\le2h.
 \label{eq:directional-bound}
\end{equation}
Indeed, on the open unit square, in the sense of distributions,
\[
 (\partial_u-\partial_v)f
 =\tfrac12\partial_u\partial_v(A-C),
\]
because $\partial_u A=-2f$ and $\partial_v C=-2f$.  The mixed derivative
of the uniform averaging kernel is the signed measure at its four corners
with total variation $h^{-2}$.  Convolution on $Q_h$ therefore gives a bound
$\|A-C\|_1/(2h^2)\le2h$ by \eqref{eq:deficit-consequences}.
Only averaging boxes contained in the original square are used here, so no
distributional boundary term from the zero extension enters this estimate.
The box average of a bounded function is Lipschitz: translating either
coordinate changes the averaging rectangle by area at most a constant
times the translation distance.  Consequently the weak directional
derivative above agrees with its almost-everywhere derivative, and the
one-dimensional fundamental theorem applies on almost every diagonal slice.

For each $s$ average $f_h(u,s-u)$ over the interval of $u$ for which
$(u,s-u)\in Q_h$, and denote the average by $j(s)$.  Put $j(s)=0$ when
that interval is empty.  This defines a measurable function with
$|j|\le1$.  On an interval $I$ of length $\ell$, an absolutely continuous
function $g$ satisfies
\[
 \int_I|g-g_I|
 \le\frac1\ell\int_I\int_I|g(x)-g(y)|\,dx\,dy
 \le\ell\int_I|g'|.
\]
The diagonal coordinates $(s,u)$ have Jacobian one and their intervals
have length at most one.  Integrating this inequality and using
\eqref{eq:directional-bound} gives
$\|f_h-j(u+v)\|_{L^1(Q_h)}\le2h$.  In particular no inverse lower bound
on the length of a short diagonal is required.  Outside $Q_h$ the area is
at most $4h$ and the pointwise difference is at most two.  Together with
\eqref{eq:box-error}, this gives
\[
 \|f-j(u+v)\|_{L^1([0,1]^2)}\le52h+2h+8h=62h.
\]
Round $j$ to a sign, with a fixed convention at zero.  For a sign $z$ and
$r\in[-1,1]$, $|z-\operatorname{sgn}(r)|\le2|z-r|$.  Thus the resulting
$f_1$ has distance at most $124h\le128h$.

For any two sign functions $f,g$, expansion of the equality indicator
gives the useful bound
\begin{equation}
 |\mathcal E(f)-\mathcal E(g)|
 \le\frac12\int_{[0,1]^2}|f-g|
       \bigl(uv+(1-u)(1-v)\bigr)\,du\,dv
 \le\|f-g\|_1.
 \label{eq:energy-lipschitz}
\end{equation}
Hence $\mathcal E(f_1)\le1/12+h^6+128h\le1/12+129h$, which is stronger
than the stated bound.  If $h\ge1/8$, any fixed diagonal sign function
suffices: the distance is at most two, the energy is at most $1/4$, and
both claimed bounds follow.  This completes the lemma for $0<h\le1$.
\end{proof}

\begin{proof}[Proof of Theorem~\ref{thm:continuous}]
Assume first that $\varepsilon>0$, put $h=\varepsilon^{1/6}$, and suppose
$\alpha=257h\le1/4096$.  Lemma~\ref{lem:diagonal-approximation} gives
$f_1(u,v)=b(u+v)$ with energy excess at most $\alpha$.
Let $t=\alpha^{1/4}\le1/8$.
By \eqref{eq:deficit-consequences} there is a row at height
$v\in[t,2t]$ whose one-dimensional deficit is at most $16\alpha/t$.
Its sign profile is $b$ restricted to $[v,v+1]$.  By
Lemma~\ref{lem:one-dimensional}, it has a threshold approximation of
error at most $4(16\alpha/t)^{1/3}<11\alpha^{1/4}$.
Restrict that threshold to $[v,1]$ and extend it to $[0,1]$ as a threshold.
The additional interval has length at most $2t$ and costs at most $4t$.
Thus $b$ on $[0,1]$ is within $15\alpha^{1/4}$ of a threshold $b_0$.
Selecting a row with $v\in[1-2t,1-t]$ in the same way gives a threshold
$b_1$ within $15\alpha^{1/4}$ of $s\mapsto b(1+s)$ on $[0,1]$.

Let $\widetilde b(s)=b_0(s)$ for $s<1$ and
$\widetilde b(s)=b_1(s-1)$ for $s\ge1$, and put
$f_2(u,v)=\widetilde b(u+v)$.  The density of $u+v$ on the square is
$s$ for $0\le s\le1$ and $2-s$ for $1\le s\le2$, so it is at most one.
Consequently
\begin{equation}
 \|f_1-f_2\|_1\le\zeta:=30\alpha^{1/4}.
 \label{eq:two-half-error}
\end{equation}
Each coordinate operator $\mathcal T$ has $L^2$ operator norm at most
one, by Cauchy--Schwarz on each unit interval and Fubini.  Moreover,
$\|f_1-f_2\|_2^2=2\|f_1-f_2\|_1$.  Writing
$\mathcal D=\mathcal T_u-\mathcal T_v$, the triangle inequality followed
by $(x+y)^2\le2x^2+2y^2$ gives
\begin{equation}
 \|\mathcal D f_2\|_2^2
 \le32\alpha+16\zeta\le512\alpha^{1/4}.
 \label{eq:two-half-operator}
\end{equation}

Set $m_0=|\{s\in[0,1]:b_0(s)=b_1(s)\}|$ and
$F(t)=\int_0^t(b_0(s)+b_1(s))\,ds$.  Directly integrating the two
coordinate operators yields
\begin{equation}
 \mathcal D f_2(u,v)=F(v)-F(u),
 \qquad
 \|\mathcal D f_2\|_2^2=2\int_0^1(F-\overline F)^2.
 \label{eq:agreement-variance}
\end{equation}
Because $b_0,b_1$ are thresholds, the nonzero derivative of $F$ is
supported on at most two intervals, and is constantly $2$ or $-2$ on
each.  If $m_0>0$, one of those intervals has length
$\ell\ge m_0/2$.  For any constant $c$, the integral of $(F-c)^2$ on
that interval is at least
\[
 \min_d\int_0^\ell(2s-d)^2\,ds=\ell^3/3.
\]
Taking $c=\overline F$ in \eqref{eq:agreement-variance} and using
\eqref{eq:two-half-operator}, we obtain
$m_0^3\le6144\alpha^{1/4}$ and hence
\begin{equation}
 m_0\le19\alpha^{1/12}.
 \label{eq:agreement-bound}
\end{equation}

Replace $b_1$ by $-b_0$ and call the resulting diagonal function $f_3$.
Since the diagonal density is at most one,
$\|f_2-f_3\|_1\le2m_0$.  If $b_0$ is positive before $c$ and negative
after $c$, then
\[
 \int f_3=\int_0^1(2s-1)b_0(s)\,ds=-2c(1-c).
\]
Reversing the threshold orientation reverses the mean.  The distances
from $f_3$ to the two canonical phases are $2c$ and $2(1-c)$, in some
order, since the density weights of the two halves add to one.
Therefore its distance to a canonical phase is at most
$2|\int f_3|$.  With $d=\|f-f_2\|_1$, the triangle inequality now gives
\[
 \min_\sigma\|f-\sigma f_*\|_1
 \le3d+6m_0+2\left|\int f\right|.
\]
Using $d\le128h+30\alpha^{1/4}$,
\eqref{eq:agreement-bound}, and $|\int f|\le3h^3$, the right-hand side
is at most
\[
 384h+90\alpha^{1/4}+114\alpha^{1/12}+6h^3
 \le594\alpha^{1/12}.
\]
Here $0<h\le1$, $\alpha=257h\le1$, and each smaller power term is
bounded by the displayed coefficient times $\alpha^{1/12}$.
Since $257^{1/12}\le2$, this is at most
$1188h^{1/12}\le4096\varepsilon^{1/72}$.

If $257h>1/4096$, the distance between any two sign functions is at most
two.  The inequality $h^{1/12}\ge1/4$ at the cutoff, and monotonicity
thereafter, give the asserted estimate immediately.  Finally, if
$\varepsilon=0$, apply the positive-excess result for every positive
excess parameter.  The minimum of the distances to the two fixed phases
must then be zero.  One of those two phases attains that minimum, proving
the zero-excess case without a choice of a limiting phase.
\end{proof}

\subsection{Finite rectangular grids with unary costs}

For integers $m,n\ge1$, let $J_{m,n}$ have vertex set
$\{0,\ldots,m-1\}\times\{0,\ldots,n-1\}$, where two vertices are adjacent
exactly when both coordinate orders strictly agree.  In this section write
\[
 V=mn,\qquad \rho=\frac1m+\frac1n,\qquad
 u_i=\frac{i+1/2}{m},\qquad v_j=\frac{j+1/2}{n}.
\]
Thus $V$ is the number of vertices of the comparison graph, not the
number of points of the eventual geometric drawing.  Define
\[
 c^+_{ij}=\mathbf1\{u_i+v_j<1\},\qquad c^-_{ij}=1-c^+_{ij},
 \qquad s^\sigma_{ij}=2c^\sigma_{ij}-1.
\]
In particular a diagonal cell centre has spin $-1$ in the positive phase.

Let $U_x:\{0,1\}\to\mathbb R$ be unary costs, and suppose
\[
 |U_x(1)-U_x(0)|\le\alpha V\quad\hbox{for all }x,
 \qquad \alpha\ge0.
\]
No bound is imposed on their additive constants.  Write
\[
 E(c)=\Phi_J(c)+\sum_xU_x(c_x),
\]
where $\Phi_J$ counts monochromatic edges of $J_{m,n}$.
For a set $F$ of grid vertices, $c\oplus F$ denotes the colouring obtained
by flipping exactly $F$.  Define the exact finite quantities
\begin{equation}
 d_x^\sigma=E(c^\sigma\oplus\{x\})-E(c^\sigma),\qquad
 L_\sigma=E(c^\sigma)+\sum_x\min(0,d_x^\sigma),\qquad
 L=\min(L_+,L_-).
 \label{eq:finite-linear-floor}
\end{equation}
For $w\ge0$ let
\[
 D_w=\{(i,j):|u_i+v_j-1|\le w\},\qquad
 C_w=\{(i,j):u_i+v_j\le w\ \hbox{or}\ u_i+v_j\ge2-w\}.
\]

\begin{theorem}[Finite-grid localization and energy bounds]
\label{thm:grid}
Let $a$ satisfy $\alpha+\rho\le a\le2^{-4096}$.
For every rectangular aspect ratio the following statements hold.
\begin{enumerate}
 \item Every minimizer $c$ of $E$ has a phase $\sigma$ such that, with
 $F=\{x:c_x\ne c_x^\sigma\}$,
 \[
 |F|\le52aV,\qquad F\subseteq D_{220a}\cup C_{220a}.
 \]
 \item The minimum energy belongs to
 \[
 [L-2^{32}a^3V^2,\ L+2^{32}a^3V^2].
 \]
 For either prescribed phase, the explicit set
 $F_\sigma^0=\{x:d_x^\sigma<0\}$ has size at most $52aV$ and
 \[
 E(c^\sigma\oplus F_\sigma^0)\le L_\sigma+2^{32}a^3V^2.
 \]
 \item For every $F\subseteq D_{220a}\cup C_{220a}$,
 \[
 E(c^\sigma\oplus F)\ge L_\sigma-2^{32}a^3V^2.
 \]
 This last conclusion requires no optimality of $c^\sigma\oplus F$.
 In fact it holds for arbitrary real unary costs, without their oscillation
 bound, whenever $\rho\le a\le1/65536$.
\end{enumerate}
For rational unary costs, both $L_\sigma$ and the two prescribed-phase
colourings are computable by polynomially many exact arithmetic operations.
\end{theorem}

\begin{proof}
\emph{From cells to the continuous objective.}
Extend a grid colouring to the sign function $f_c$ constant on its
$1/m$ by $1/n$ cells.  Comparable pairs in distinct rows and columns
contribute exactly $\Phi_J(c)/V^2$.  All other contributions are
nonnegative and their total uncoloured measure is
\[
 \frac14-\frac{\binom m2\binom n2}{m^2n^2}
 =\frac1{4m}+\frac1{4n}-\frac1{4mn}\le\frac\rho4.
\]
We will use the weaker bound
\begin{equation}
 0\le\mathcal E(f_c)-\Phi_J(c)/V^2\le\rho/2.
 \label{eq:cell-comparison}
\end{equation}
The functions $f_{c^+}$ and $f_*$ can disagree only where
$|u+v-1|\le\rho/2$.  This band has area at most $\rho$, whence
$\|f_{c^+}-f_*\|_1\le2\rho$, and the analogous assertion holds for
the negative phase.  By \eqref{eq:energy-lipschitz},
$\mathcal E(f_{c^+})\le1/12+2\rho$.
For a minimizer $c$, comparison with $c^+$ gives
$\Phi_J(c)\le\Phi_J(c^+)+\alpha V^2$.  Therefore
\[
 \mathcal E(f_c)\le1/12+\alpha+\tfrac52\rho\le1/12+5a.
\]
Theorem~\ref{thm:continuous} and the canonical cell discrepancy give a
phase whose disagreement fraction $\mu=|F|/V$ satisfies
\begin{equation}
 \mu\le2048(5a)^{1/72}+\rho\le2^{-20}.
 \label{eq:initial-fraction}
\end{equation}
For the final numerical bound, use $5<8$, $a\le2^{-4096}$, and
$\rho\le a$; the first term is less than $2^{11-4093/72}<2^{-45}$.
This initial estimate follows from the minimum condition.

\emph{Canonical local fields.}
For $p=(u,v)$ let $Q_p$ be the union of the lower-left and upper-right
rectangles comparable to $p$.  With $s=1-u-v$, direct integration gives
\[
 \int_{Q_p}f_*=-s(1-|s|).
\]
For $s\ge0$, the lower-left rectangle contributes $uv$, and the
upper-right rectangle contributes $s^2-(1-u)(1-v)$; reflection gives the
other case.  The union of cells corresponding to the finite neighbours of
a cell centre differs from $Q_p$ only in its row and column strips,
of total area at most $\rho$.  The sign-function discrepancy contributes
at most another $2\rho$.  A single grid flip changes the cost by minus
the vertex spin times its neighbour-spin sum.  Consequently, in either
canonical phase,
\begin{equation}
 \left|\frac{d_x^{\rm grid}}V-\ell(u_i,v_j)\right|\le6\rho,
 \qquad
 \ell(u,v)=|1-u-v|(1-|1-u-v|).
 \label{eq:canonical-field}
\end{equation}
At a diagonal centre the continuous field is zero, so the strict tie
convention causes no change to this bound.  Including the unary cost gives
\begin{equation}
 d_x^\sigma\ge V\ell(u_i,v_j)-6aV.
 \label{eq:field-with-unaries}
\end{equation}
If $t\in[0,1]$ and $t(1-t)\le z$, then either $t\le2z$ or
$1-t\le2z$, by considering $t\le1/2$ and $t\ge1/2$ separately.
Thus a small value of $\ell$ places the vertex in a narrow diagonal strip
or in the two corner regions.

\emph{Finite strip and corner counts.}
A rectangle with side lengths $r,t$ contains at most
\[
 (mr+1)(nt+1)\le V(r+\rho)(t+\rho)
\]
cell centres.  For $\rho\le w\le1$, counting by rows gives
$|D_w|\le(2w+\rho)V\le3wV$, and bounding each corner triangle by
a square gives $|C_w|\le8w^2V$.  Two comparable points of $D_w$, ordered
coordinatewise, have the sum of their coordinate differences at most
$2w$.  For any strip point its strip neighbours therefore lie in two
squares of side $2w$, giving at most $18w^2V$ possibilities.  Hence for
every $F\subseteq D_w\cup C_w$,
\begin{equation}
 \deg_F(x)\le26w^2V\quad(x\in D_w),\qquad
 |E(J_{m,n}[F])|\le83w^3V^2.
 \label{eq:strip-counts}
\end{equation}
For the edge bound, strip--strip, strip--corner, and corner--corner
edges contribute at most $27w^3V^2$, $24w^3V^2$, and $32w^4V^2$,
respectively.  Counting an edge more than once when regions overlap only
increases these upper bounds.  The degree assertion is made only for strip
vertices: corner vertices need not have degree of order $w^2V$.

\emph{Bootstrap from initial stability.}
For a disagreeing vertex $x\in F$, changing its colour back in the
minimizer has cost change
\[
 -d_x^\sigma-2\sum_{y\in F:\,xy\in E(J_{m,n})}s_x^\sigma s_y^\sigma
 \ge0.
\]
Thus $d_x^\sigma\le2\deg_F(x)$, and
\eqref{eq:field-with-unaries} yields
\begin{equation}
 \ell(u_i,v_j)\le6a+2\deg_F(x)/V\qquad(x\in F).
 \label{eq:local-optimality-profile}
\end{equation}
Initially $\deg_F(x)\le\mu V$.  The observation following
\eqref{eq:field-with-unaries} gives
\[
 F\subseteq D_w\cup C_w,
 \qquad w=4\mu+12a.
\]
By \eqref{eq:initial-fraction}, $\rho\le w<1/4$.  For
$x\in F\cap D_w$, \eqref{eq:strip-counts} improves
\eqref{eq:local-optimality-profile} to
$\ell\le6a+52w^2$.  In $D_w$ we have
$\ell\ge|1-u-v|/2$, so
$F\cap D_w\subseteq D_{12a+104w^2}$.  The direct row count and the
corner count now imply
\[
 \mu\le25a+216w^2
 \le25a+6912\mu^2+62208a^2.
\]
Since $\mu\le2^{-20}$ and $a\le2^{-4096}$, we have
$6912\mu^2\le\mu/2$ and $62208a^2\le a$.
Absorbing these terms proves $\mu\le52a$.
Substitution in the definition of $w$ gives $w\le220a$, establishing
the first conclusion.  No small degree bound for corner vertices has
been used in this absorption.

\emph{Exact correction for simultaneous flips.}
For every flip set $F$ and arbitrary real unary costs,
\begin{equation}
 E(c^\sigma\oplus F)-E(c^\sigma)
 =\sum_{x\in F}d_x^\sigma
  +2\sum_{\{x,y\}\in E(J_{m,n}[F])}s_x^\sigma s_y^\sigma.
 \label{eq:simultaneous-flips}
\end{equation}
Indeed, an edge with both endpoints flipped keeps its equality status,
while its two single-flip contributions sum to
$-2s_x^\sigma s_y^\sigma$.  The second sum restores exactly this
discrepancy; unary terms already add independently.

If $F\subseteq D_{220a}\cup C_{220a}$ and
$\rho\le a\le1/65536$, \eqref{eq:strip-counts} bounds the absolute
correction by
\[
 2\cdot83\cdot220^3a^3V^2<2^{32}a^3V^2.
\]
The linear sum is at least $\sum_x\min(0,d_x^\sigma)$.
This proves the third conclusion, including its stronger independence
from the unary oscillation and from optimality.  Applying it to the
localized minimizer proves the lower bound in the second conclusion.

For the upper bound in a prescribed phase, take exactly
$F_\sigma^0=\{x:d_x^\sigma<0\}$.  By
\eqref{eq:field-with-unaries}, this set is contained in
$D_{12a}\cup C_{12a}$, and its cardinality is at most
\[
 (36a+1152a^2)V\le52aV.
\]
Its linear sum in \eqref{eq:simultaneous-flips} is the defining sum for
$L_\sigma$, while its correction is bounded by
$2\cdot83\cdot12^3a^3V^2<2^{32}a^3V^2$.
Choosing the phase with smaller $L_\sigma$ proves the upper bound for
the minimum.  Finally, all quantities in \eqref{eq:finite-linear-floor}
are finite sums and single-vertex comparisons.  They require no search
over colourings and can be evaluated exactly in polynomially many
operations for rational unary costs.
\end{proof}

\subsection{Freezing probes with the grid colouring fixed}

The following lemma is about an arbitrary finite simple graph.  In the
application its vertices are geometric edges and its adjacency relation
is proper crossing.  Thus its conclusion includes interactions between
the probes themselves.

\begin{proposition}[Probe freezing]
\label{prop:freezing}
Let a finite simple graph have vertex partition $X,T,Z$, with induced
subgraph on $X$ equal to $J_{m,n}$.  The objective is the number of
monochromatic graph edges, and all vertices may be coloured freely.
Use $V=mn$ and $\rho=1/m+1/n$.  Suppose every vertex of $X$ has at most
$\alpha V$ neighbours outside $X$, every vertex of $T$ has at most
$\beta V$ neighbours outside $X$, and
\[
 \alpha,\beta\ge0,\qquad \alpha+\rho\le a\le2^{-4096}.
\]
For $p\in T$ and $\sigma\in\{+,-\}$ let
\[
 h_p^\sigma=\sum_{x\in N(p)\cap X}s_x^\sigma.
\]
This is its colour-1 cost minus its colour-0 cost when the grid is
canonical.  Choose $y_p^+=1$ when $h_p^+<0$ and $y_p^+=0$ otherwise,
and put $y_p^-=1-y_p^+$.  These are best responses in both phases,
including their complemented zero-field tie convention.

Put $b=104a+\beta$, and suppose that $Q\ge0$ satisfies
\begin{equation}
 |\{p\in T:|h_p^\sigma|\le bV\}|
 \le Q(b+\rho)V
 \quad\hbox{for both phases}.
 \label{eq:soft-probe-count}
\end{equation}
Then every optimal colouring $\chi$ has a phase $\sigma$ and a colouring
$\psi$ such that only vertices of $T$ are changed, every probe is coloured
$y_p^\sigma$, and
\begin{equation}
 0\le\Phi(\psi)-\Phi(\chi)
 \le E_{\rm fr}:=2Q(b+\rho)bV^2.
 \label{eq:freezing-error}
\end{equation}
The unchanged grid differs from $c^\sigma$ on at most $52aV$ vertices,
all in $D_{220a}\cup C_{220a}$.  If $\OPT_{\rm fr}$ is the minimum over
colourings with the probes in either of the two frozen phases, then
\[
 0\le\OPT_{\rm fr}-\OPT\le E_{\rm fr}.
\]
It suffices for \eqref{eq:soft-probe-count} to have the stronger estimate
$|\{p:|h_p^\sigma|\le tV\}|\le Q(t+\rho)V$ for
$0\le t\le t_0$, together with $b\le t_0$.
\end{proposition}

\begin{proof}
Fix the colours outside $X$ in a global optimum.  Its restriction
to $X$ minimizes the resulting grid objective.  Each conditional unary
oscillation is at most the number of outside neighbours, hence at most
$\alpha V$.  Theorem~\ref{thm:grid} gives a phase $\sigma$ and a
disagreement set $F$ with $|F|\le52aV$ and the stated support.  Keep this
grid colouring fixed for the rest of the proof.

For every probe, its field from $X$ differs from $h_p^\sigma$ by at most
$2|F|\le104aV$.  Under any colouring of all other vertices, its field
from neighbours outside $X$ has absolute value at most $\beta V$.
Thus the full field differs from $h_p^\sigma$ by at most $bV$.
If $|h_p^\sigma|>bV$, the sign of the full field is the same as the
canonical sign, whatever the outside colours are.  Single-vertex
optimality of $\chi$ therefore forces this hard probe already to have
colour $y_p^\sigma$.

Recolour each remaining, soft probe to $y_p^\sigma$, in any order.
At every step the grid remains fixed.  The outside field still has
absolute value at most $\beta V$ even after earlier probe changes, since
all probe--probe neighbours were included in that degree bound.  The
cost increase for the current recolouring is at most the magnitude of
the current full field, which is bounded by
\[
 |h_p^\sigma|+104aV+\beta V\le2bV.
\]
There are at most $Q(b+\rho)V$ soft probes by
\eqref{eq:soft-probe-count}.  Summing the sequential cost changes gives
the upper bound in \eqref{eq:freezing-error}.  The lower bound follows
because $\chi$ is a global optimum and $\psi$ is a feasible colouring.
Only probes were changed, so the same grid and its localization remain.

The frozen feasible set is a subset of the original feasible set and
contains the constructed $\psi$.  Its optimum is therefore between
$\Phi(\chi)$ and $\Phi(\psi)$, proving the final optimum bound.
\end{proof}

The last clause of Theorem~\ref{thm:grid} can now be applied to this
unchanged grid under the newly frozen unary costs.  It applies because
the support is localized, whether or not this grid minimizes the new
conditional objective.  This is the reason for retaining the arbitrary
localized-colouring estimate in addition to the optimum bracket.
